% !TEX program = LuaLaTeX+se
% !TEX encoding = UTF-8 Unicode

\documentclass[12pt, a4paper]{article}
\usepackage[latin, polish]{babel} % ustawienia języka
\usepackage[T1]{fontenc} % przełączenie na fonty EC (256 znaków)
\usepackage[utf8]{luainputenc} % kodowanie
\usepackage{fullpage} % redukcja marginesów
\usepackage{setspace} % odstępy między liniami
\usepackage{geometry} 
\geometry{
	left=12mm,
	right=12mm,
	top=12mm,
	bottom=12mm,
	marginparwidth=0mm,
	marginparsep=0mm
	}

\usepackage{fontspec}
\usepackage{libertine} % ust. czcionki
\pagestyle{empty}	% usuwanie nagłówków/stopek z dokumentu
\usepackage[indent=0pt]{parskip} % ustawianie paragrafów
\setstretch{1.5} 

%  Konfiguracja kolumn
\usepackage{paracol}
\setcolumnwidth{8.3cm,4.3cm}
\setlength{\columnsep}{0.4cm}

\usepackage{gregoriotex}
\gresetlinecolor{gregoriocolor}
\grechangestyle{initial}{\setmainfont{VectisMonumentalW01-Regular}\fontsize{23}{23}\selectfont}
% \grechangedim{spacelinestext}{0.7 cm}{scalable}
% \grechangedim{spaceabovelines}{0.43 cm}{scalable}
% \gresetheadercapture{commentary}{grecommentary}{string}
% \grecommentary{}
\gresetheadercapture{annotation}{greannotation}{string}
% Format annotation above initial
\grechangestyle{annotation}{\small\bfseries}%
%\grechangedim{intersyllablespacenotes}{0.18cm}{scalable}%
%\grechangestyle{commentary}{\it}
%\gresetmodenumbersystem{roman-majuscule}
\usepackage{titlesec} % ustawianie nagłówków
\titlespacing*{\section}{0pt}{0pt}{3mm}	
\titleformat{\section}[hang]{\filcenter\large\scshape\bfseries}{}{0pt}{} % format nagłówków

\newfontface\vectis{VectisMonumentalW01-Regular}
\def\greinitialformat#1{%
{\fontsize{24}{24}\vectis #1}}

\titlespacing*{\section}
{0pt}{0pt}{0pt}

\begin{document}
\large
\begin{spacing}{0.7}
\begin{center}
	\section*{Jesu corona celsior (Conf. non Pont.)}
	\small\textit{Antiphonale Monasticum, Solesmes, 1934, p. 672}
	\vspace{12mm}
\end{center}	
\end{spacing}
\greannotation{Hymn.} 
\greannotation{8.} 
\gabcsnippet{(c4)JE(g)su,(g') co(f)ró(e)na(f) cél(g)si(fg)or,(g_') (,)
Et(h) vé(jk)ri(j)tas(j) sub(g)lí(h')mi(g)or,(f.) (;)
Qui(f) con(g')fi(f)tén(g)ti(h) sér(j')vu(i)lo,(giHG'_) (,)
Red(g)dis(gh) per(f)én(e')ne(f) prǽ(g')mi(h)um.(g.) (::)

2. Da(g) sup(g')pli(f)cán(e)ti(f) cœ(g)́tu(fg)i,(g_') (,)
Ob(h)tén(jk)tu(j) hu(j)ius(g) óp(h')ti(g)mi,(f.) (;)
Re(f)mis(g')si(f)ó(g)nem(h) crí(j')mi(i)num,(giHG'_) (,)
Rum(g)pén(gh)do(f) ne(e')xum(f) vín(g')cu(h)li.(g.) (::)

3. An(g)ni(g') re(f)cúr(e)so(f) tém(g)po(fg)re,(g_') (,)
Di(h)es(jk) il(j)lú(j)xit(g) lú(h')mi(g)ne,(f.) (;)
Quo(f) Sanc(g')tus(f) hic(g) de(h) cór(j')po(i)re(giHG'_) (,)
Po(g)lum(gh) mi(f)grá(e')vit(f) prǽ(g')pot(h)ens.(g.) (::)

4. Hic(g) va(g')na(f) ter(e)ræ(f) gáu(g)di(fg)a,(g_') (,)
Et(h) lu(jk)cu(j)lén(j)ta(g) prǽ(h')di(g)a,(f.) (;)
Pol(f)lú(g')ta(f) sor(g)de(h) dé(j')pu(i)tans,(giHG'_) (,)
O(g)vans(gh) te(f)net(e') cæ(f)lés(g')ti(h)a.(g.) (::)

5. Te,(g) Chris(g')te,(f) Rex(e) pi(f)ís(g)si(fg)me,(g_') (,)
Hic(h) con(jk)fi(j)tén(j)do(g) jú(h')gi(g)ter,(f.) (;)
Cal(f)cá(g')vit(f) hos(g)tem(h) fór(j')ti(i)ter,(giHG'_) (,)
Su(g)pér(gh)bum(f) ac(e') sa(f)tél(g')li(h)tem.(g.) (::)

6. Vir(g)tú(g')te(f) cla(e)rus(f) et(g) fi(fg)de,(g_') (,)
Con(h)fes(jk)si(j)ó(j)nis(g) ór(h')di(g)ne,(f.) (;)
Je(f)jú(g')na(f) mem(g)bra(h) dé(j')fe(i)rens,(giHG'_) (,)
Da(g)pes(gh) su(f)pér(e')nas(f) ób(g')ti(h)net.(g.) (::)

7. Pro(g)ín(g')de(f) te,(e) pi(f)ís(g)si(fg)me,(g_') (,)
Pre(h)cá(jk)mur(j) om(j)nes(g) súp(h')pli(g)ces,(f.) (;)
Ut(f) hu(g')jus(f) al(g)mi(h) grá(j')ti(i)a(giHG'_) (,)
No(g)bis(gh) re(f)mít(e')tas(f) dé(g')bi(h)ta.(g.) (::)

8. Gló(g)ri(g')a(f) Pa(e)tri(f) Dó(g)mi(fg)no,(g_') (,)
Gló(h)ri(jk)a(j) U(j)ni(g)gé(h')ni(g)to,(f.) (;)
U(f)na(g') cum(f) Sanc(g)to(h) Spí(j')ri(i)tu,(giHG'_) (,)
In(g) sem(gh)pi(f)tér(e')na(f) sǽ(g')cu(h)la.(g.) (::)
A(ghg)men.(fg..) (::)}
\vspace{4mm}
\setstretch{1}% odstęp między liniami
1. Jezu, korono bezcenna! Prawdo, w swem słowie niezmienna! Co swego sługę wiernego Domieszczasz szczęścia wiecznego.\\
2. Racz trzodkę, Ciebie modlącą, Przez Jego prośbę znaczącą, Z win oswobodzić więzienia, I żądz cielesnych brzmienia.\\
3. Po rocznym czasie przegnanym, Znów się dniem cieszym żądanym, W którym, z tej ziemi, nasz święty Został do nieba przyjęty.\\
4. On zmienne świata wygody, I znakomite zagrody, Poznał, iż brudzą niecnotą: I szedł do nieba z ochotą.\\
5. Tu się Cię, Chryste! trzymając, Tu Cię swym Królem wyznając, Podeptał sidła szatana - Srogiego piekieł tyrana.\\
6. Z męztwa i z wiary swej znany, Czci Pana szczerze oddany: Ciału nie dawał swobody, I poszedł w Niebo na gody.\\
7. Stąd Ciebie, Panie litośny! Prosi nasz poczet żałośny: Przez Jego ważne staranie, Chciej nam darować karanie.\\
8. Cześć, chwała, Ojcu wiecznemu, Jego Synowi jednemu, I od obojga Duchowi, Na wieki niechaj się słowi.\footnote{tłum. Hymny ojców świętych i wielkich w Kościele mężów z brewiarza i mszału, przeł. X.A.Kitkiewicz, Wilno, 1848}\\
Amen
\end{document}